% =========================================================
% DIAGRAMA: Secuencia ISO-TP multi-frame (FF → FC → CF×N)
% Archivo imagen: Ilustraciones/isotp_multiframe.png
%
% QUÉ DIBUJAR EN DRAW.IO:
%   Diagrama de secuencia (sequence diagram) con dos actores verticales:
%   Izquierda: "Escáner / Raspberry Pi"   Derecha: "ECU"
%
%   Flujo (flechas horizontales de izq a der y der a izq):
%
%   Escáner ──────[ Request SF: 02 09 02 ]──────────────► ECU
%           (petición VIN, Single Frame, Mode 09 PID 02)
%
%   ECU     ──────[ FF: 10 14 49 02 01 57 30 30 ]───────► Escáner
%           (First Frame, longitud total=20 bytes, primeros 6 bytes payload)
%
%   Escáner ──────[ FC: 30 00 00 AA AA AA AA AA ]───────► ECU
%           (Flow Control: ContinueToSend, BS=0, STmin=0)
%
%   ECU     ──────[ CF SN=1: 21 30 30 30 30 30 30 30 ]──► Escáner
%   ECU     ──────[ CF SN=2: 22 30 30 30 30 30 30 30 ]──► Escáner
%           (Consecutive Frames hasta completar 20 bytes)
%
%   - Línea de vida vertical punteada para cada actor
%   - Cada flecha etiquetada con el tipo de trama y bytes hex
%   - Nota al lado del FF: "Payload total: 20 bytes (3 header + 17 VIN)"
%   - Nota al lado del FC: "BS=0: sin límite de bloques"
%   - Franja temporal destacando "Segmento multi-frame"
%   - Título: "Secuencia ISO-TP: petición VIN (20 bytes)"
% =========================================================

\begin{figure}[!htbp]
\centering
\includegraphics[width=0.85\textwidth]{Ilustraciones/isotp_multiframe.png}
\caption{Secuencia de intercambio ISO-TP para la obtención del VIN (20 bytes). El escáner inicia con una petición Single Frame; la ECU responde con First Frame seguido de Consecutive Frames tras recibir el Flow Control.}
\label{fig:isotp_multiframe}
\end{figure}
